% =====================================================================
%  Phase C --- The MQTT contract
% =====================================================================
\section{The Transport Contract}
\label{sec:phaseC_protocol}

\subsection{Why MQTT and not HTTP}

The Cloud must be able to reach the robot, not only the other way round.
With HTTP the robot would have to poll for work; with MQTT it subscribes
once and the broker pushes. The broker is also the only endpoint either
side configures: neither has to know the other's address, and either can
restart without the other noticing. That is the property that makes
``the robot is a node on a network'' an architectural fact rather than a
turn of phrase.

It is additionally the protocol the rest of the agricultural-IoT world
speaks, so the same robot could report to a different back end without a
line changing in \filename{agri/protocol.py}.

\subsection{Topics, namespaced by node}

\begin{table}[!t]
\caption{The topic contract. Robot-to-Cloud topics are namespaced per
node; the Cloud subscribes with MQTT wildcards.}
\label{tab:topics}
\centering\footnotesize
\begin{tabular}{@{}llp{0.30\columnwidth}@{}}
\toprule
topic & direction & payload \\
\midrule
\topicname{agri/v1/request}              & Cloud $\rightarrow$ all & signed, readable \\
\topicname{agri/v1/ack/<node>}           & node $\rightarrow$ Cloud & plain progress \\
\topicname{agri/v1/report/<node>}        & node $\rightarrow$ Cloud & sealed \\
\topicname{agri/v1/status/<node>}        & node $\rightarrow$ Cloud & plain, retained \\
\bottomrule
\end{tabular}
\end{table}

Requests stay on a \emph{flat} topic because a request addresses the
fleet, not a specific node. Everything a node says is namespaced by its
own identifier, and the Cloud subscribes with \texttt{+} to hear all of
them, keeping nodes in a dictionary rather than a single variable. There
is one mobile node today; the layout is what stops adding a second from
being a protocol change.

\subsection{Node kinds}

Every status message carries a \param{node\_kind} field:
\texttt{"mobile"} for a robot that drives between stations,
\texttt{"fixed"} for a stationary microcontroller bolted to a post. The
field is carried now, with one mobile node and no fixed ones, so that
adding fixed sensors later does not change the protocol --- and so that
the priority rule it exists for (prefer the mobile reading when a robot
is parked at a station a fixed node also covers) has somewhere to live.

\subsection{Three choices that were not defaults}

\paragraph{QoS 1, not 0 or 2} Zero would drop a request while the robot
is mid-reconnect. Two costs a four-way handshake for no benefit here,
because the request identifier already makes a repeat harmless --- and
the store deduplicates on (label, timestamp, robot), which catches
exactly the case QoS~1 permits and nothing else.

\paragraph{Status is retained; nothing else is} A dashboard opened at
any moment should immediately know whether the robot is there, and MQTT's
retained flag gives precisely that. Retaining \emph{reports} instead
would mean every reconnecting client re-receives the last measurement and
files it again, which is how duplicates appear in a store that has no
reason to have any.

\paragraph{A last will} If the robot process dies, the broker tells the
Cloud on the node's own status topic. A dashboard that shows a robot as
online because nobody said otherwise is worse than one that shows
nothing.

\subsection{Two failures the transport layer had to survive}

Both were observed, and both are now pinned by the suite.

\paragraph{A broker that is not up yet} The robot node originally exited
when it could not connect, which made the launch order of two terminals
load-bearing --- start the simulation before the broker and the robot was
simply gone, with the failure scrolled off the top of a 140-line startup
log. It now waits, retries with a bounded backoff, re-subscribes on
\emph{every} connect rather than only the first, says so when the broker
goes away, and complains on the wall clock if none ever appears.

\paragraph{An exception inside an MQTT callback} \texttt{paho} does not
guard its callbacks: an exception raised inside one kills its network
thread. The observed symptom was a robot that looked perfectly healthy
--- its status timer kept publishing, because that runs on a different
thread --- while every request the Cloud sent was delivered to a client
that was no longer reading. Forty minutes of a robot that would not move
and a log that said nothing.

The cause was the status callback asking the driver for a pose before
the simulator had published any odometry. Two rules came out of it and
are asserted statically: no MQTT callback may let an exception escape,
and \texttt{status()} must never raise even with no odometry at all --- a
status with no position is still worth sending, because ``I am here and
I am connected'' is exactly what the Cloud needs at that instant.
